% Options for packages loaded elsewhere
\PassOptionsToPackage{unicode}{hyperref}
\PassOptionsToPackage{hyphens}{url}
%
\documentclass[
]{article}
\usepackage{amsmath,amssymb}
\usepackage{lmodern}
\usepackage{iftex}
\ifPDFTeX
  \usepackage[T1]{fontenc}
  \usepackage[utf8]{inputenc}
  \usepackage{textcomp} % provide euro and other symbols
\else % if luatex or xetex
  \usepackage{unicode-math}
  \defaultfontfeatures{Scale=MatchLowercase}
  \defaultfontfeatures[\rmfamily]{Ligatures=TeX,Scale=1}
\fi
% Use upquote if available, for straight quotes in verbatim environments
\IfFileExists{upquote.sty}{\usepackage{upquote}}{}
\IfFileExists{microtype.sty}{% use microtype if available
  \usepackage[]{microtype}
  \UseMicrotypeSet[protrusion]{basicmath} % disable protrusion for tt fonts
}{}
\makeatletter
\@ifundefined{KOMAClassName}{% if non-KOMA class
  \IfFileExists{parskip.sty}{%
    \usepackage{parskip}
  }{% else
    \setlength{\parindent}{0pt}
    \setlength{\parskip}{6pt plus 2pt minus 1pt}}
}{% if KOMA class
  \KOMAoptions{parskip=half}}
\makeatother
\usepackage{xcolor}
\setlength{\emergencystretch}{3em} % prevent overfull lines
\providecommand{\tightlist}{%
  \setlength{\itemsep}{0pt}\setlength{\parskip}{0pt}}
\setcounter{secnumdepth}{-\maxdimen} % remove section numbering
\ifLuaTeX
  \usepackage{selnolig}  % disable illegal ligatures
\fi
\IfFileExists{bookmark.sty}{\usepackage{bookmark}}{\usepackage{hyperref}}
\IfFileExists{xurl.sty}{\usepackage{xurl}}{} % add URL line breaks if available
\urlstyle{same} % disable monospaced font for URLs
\hypersetup{
  hidelinks,
  pdfcreator={LaTeX via pandoc}}

\author{Dietmar Gerald Schrausser}
\date{09.05.2026}


\begin{document}

\hypertarget{foliumblat-iiiiv}{%
\section{Foliu{[}m{]}/Blat IIIIv}\label{foliumblat-iiiiv}}

Eine zeitgenössische deutsche Übersetzung wird präsentiert, die auf
einem Vergleich der bestehenden lateinischen, frühneuhochdeutschen
(ENHG) und englischen (Übersetzungs-) Texte basiert, um ein besseres
Verständnis zu ermöglichen. Insbesondere da der ENHG-Text im Vergleich
zum modernen Hochdeutschen schwer verständlich ist, da dieser fundierte
Deutschkenntnisse (als Muttersprache) sowie Kenntnisse verschiedener
Dialekte voraussetzt.

\hypertarget{vom-werk-des-fuxfcnften-tages}{%
\subsection{Vom Werk des fünften
Tages}\label{vom-werk-des-fuxfcnften-tages}}

\begin{quote}
\textbf{Q}Uinto die dixit deus: producant aque reptile anime viuentis.
{[}et{]} volatile super terra{[}m{]} sub firmamento celi.
Cerauitq{[}ue{]} de{[}us{]} cete grandia: {[}et{]} omnem animam viuentem
atq{[}ue{]} motabilem quas produxera{[}n{]}t aque in species suas.
{[}et{]} om{[}n{]}e volatile sec{[}un{]}d{[}u{]}m genus suum.
Uide{[}n{]}s quod esset bonum benedixit eis dicens. Crescite {[}et{]}
multiplicamini {[}et{]} replete aquas maris: auesq{[}ue{]}
multiplicentur super terram.
\end{quote}

So schmückte Gott an diesem Tag die Luft und das Wasser mit Wesen,
geflügelte gab er der Luft und schwimmende dem Wasser, die man
\emph{Kriechtiere}\footnote{`reptilia', nur in der lateinischen Fassung.}
nennt, welche mit einer gewissen Ungestümtheit umhereilen. In der Tat
findet man große Wale und Wasserlebewesen im Meer, von monströser
Gestalt, größer als Landtiere aufgrund ihrer Fülle an
Feuchtigkeit\footnote{`habundantia humoris' und `uberflüssigkeit irer
  feuchtigkeit'.}.\footnote{aus der `Historia Naturalis' (Plinius
  Secundus,
  \href{https://gallica.bnf.fr/ark:/12148/btv1b550140045/}{1250}, Buch
  9, Kapitel 2).} Man sagt, was immer in irgendeinem Teil der Natur
entsteht, existiert auch im Meer.\footnote{aus der `Historia Naturalis'
  (Plinius Secundus,
  \href{https://gallica.bnf.fr/ark:/12148/btv1b550140045/}{1250}, Buch
  9, Kapitel 1).}

Nun zu folgenden Dingen, die hinsichtlich der Entstehung der Tiere
offenkundig sind: \emph{Nach} den (i) Pflanzen kommen jene (ii) Wesen,
die fühlen und sich bewegen; obwohl selbst den Pflanzen die Pythagoreer
einen stumpfen\footnote{`stupidum sensum' und `unbrüfende
  empfindlichkeit'.} Sinn zuschreiben. Solche Lebewesen, die eindeutig
Bewegung und Empfindung aufweisen, werden von \textbf{Moses} und im
\textbf{\emph{Timaios}} in fliegende Tiere, Wassertiere und Landtiere
unterteilt.\footnote{Eine Parallele zwischen dem biblischen
  Schöpfungsbericht (Moses) und Platons Kosmologie in Timaios
  (\href{https://doi.org/10.3931/e-rara-24349}{1588}) zu ziehen, war ein
  häufiges Thema im Renaissance-Neuplatonismus oder der Scholastik.}

Kommen wir nun zu \emph{Moses}, der, nachdem er über die himmlischen
Dingen gesprochen hatte, die irdischen Tiere in angemessener Reihenfolge
erwähnte, die entweder im Wasser, auf der Erde oder in der Luft
leben.\footnote{Wahrscheinlich vom `Hexameron' (Ambrosius,
  \href{https://mdz-nbn-resolving.de/details:bsb00046506}{12th c.}).}

\begin{quote}
Si tamen inhabitare aerem volucres dici possunt.\footnote{`Falls man
  überhaupt sagen kann, dass Vögel die Luft bewohnen', Passage aus dem
  `Heptaplus' (Mirandola \& Mirandola,
  \href{https://doi.org/10.3931/e-rara-61217}{1601}, Buch I, Kapitel 5)
  hier wird das Wesen der Vögel sowie ihr Verhältnis zur \emph{Luft} als
  Lebensraum im Vergleich zu den \emph{Wassern} der biblischen
  Schöpfungsgeschichte hinterfragt.}
\end{quote}

\textbf{\emph{Belasse man diese Diskussion hier}}. Auf welche Weise
entstehen die Körper der Tiere\footnote{`corpora animalium' und `lieb
  der thier'.} aus den Elementen, was sind die basalen
Gründe\footnote{`seminarie rationes' und `besamunge{[}n{]}'.}, die Gott
in die Natur der Dinge eingepflanzt hat, wird das Leben der
unvernünftigen Tiere ebenfalls aus dem Schoß der Materie hervorgebracht,
oder geht alles Leben vielmehr von einem göttlichen Prinzip aus, wie
\textbf{Plotinus} am entschiedensten behauptet. Eine Aussage, die der
Prophet in folgender Passage zu unterstützen scheint, nachdem er
sagte:\footnote{Aus Ficinos Kommentar zu den `Enneaden' des Plotinus
  (\href{https://doi.org/10.3931/e-rara-120889}{1580}).}

\begin{quote}
producant aque reptile anime viuentis.\footnote{`Mögen die Wasser
  lebende, kriechende Geschöpfe hervorbringen'.}
\end{quote}

Später wurde hinzugefügt:

\begin{quote}
creauit deus omne anima{[}m{]} viuente{[}m{]}.\footnote{`Gott schuf jede
  lebendige Seele'.}
\end{quote}

Hierbei ist zu beachten, dass das \emph{Wasser} diese auf Gottes (i)
Geheiß hervorbrachte und Gott selbst sie dann schuf. Dort wo von Gottes
(ii) Wirken die Rede ist heisst es dann

\begin{quote}
creauit deus \emph{anima{[}m{]}} viuente{[}m{]}, \footnote{`Gott schuf
  eine lebendige Seele'.}
\end{quote}

dort hingegen, wo lediglich vom \emph{Wasser} die Rede ist, anstatt von

\begin{quote}
\emph{a{[}n{]}imam}\footnote{`einer Seele'.}
\end{quote}

von

\begin{quote}
reptile \emph{a{[}n{]}i{[}m{]}e}\footnote{`kriechende Geschöpfe'. Sowohl
  \emph{anima} als auch \emph{animus} bedeuten \emph{Seele}, beschreiben
  aber unterschiedliche Aspekte der Existenz: \textbf{\emph{Anima}}, das
  Lebensprinzip, \textbf{\emph{Animus}}, den rationalen Verstand, vgl.
  C. G. Jungs `inneres Selbst' (s. Jung,
  \href{https://books.google.com/books?id=cX3VEAAAQBAJ}{2023}).}
\end{quote}

also einem \emph{Vehikel} für die vom Wasser bereitgestellte
Seele.\footnote{Dieser Text ist eine Passage aus `Genesis ad litteram'
  (Augustinus,
  \href{https://mdz-nbn-resolving.de/details:bsb00046071}{1961}, Buch 3,
  Kapitel 13). Er behandelt die sprachlichen und theologischen Nuancen
  der Schöpfung des Wasserlebens in Genesis 1,20-21 (vgl. Jiménez de
  Cisneros, \href{https://doi.org/10.3931/e-rara-46695}{1517}).}

Wie Moses im folgenden Schöpfungstag \emph{drei} Tierarten auf der Erde
erwähnt, befinden sich die zahlreichsten und größten Tiere im Meer der
Inder, darunter sind \emph{große Fische}\footnote{`balene'.} in der
Größe von vier \emph{Iugera}\footnote{1 Iugerum = 0.25 Hektar, wohl eher
  ein Fischschwarm.}, auch sieht man \emph{Wale}\footnote{`bellue' und
  `wu{[}n{]}der thier'.} im Meer dort, wo die Sonne ihren Lauf
nimmt\footnote{Den Wendekreisen.}. Hier werden sie von gewaltigen
Wellenbergen aus den Tiefen des Meeres emporgetragen, bis sie den
Menschen sichtbar werden. Dann toben Wirbelwinde, kommt der Regen,
daraufhin Stürme, herabgestürzt von den Bergkämmen, wirbeln die Meere
vom Grund auf und reißen Meerestiere mit den Wellen aus der Tiefe. Auch
unter den Vögeln sind die größten, beinahe die größten aller Tiere, die
afrikanischen oder äthiopischen Strauße\footnote{`generis strucio
  cameli'.}, die einen Reiter auf einem Pferd überragen und an
Geschwindigkeit übertreffen.\footnote{aus der `Historia Naturalis'
  (Plinius Secundus,
  \href{https://gallica.bnf.fr/ark:/12148/btv1b550140045/}{1250}, Buch
  9, 10).} Noch viele weitere wundebare Dinge über die Natur der Vögel
und Fische werden täglich an vielen verschiedenen Orten
erlebt\footnote{`Multo mirabilius de naturis auium {[}et{]} piscium
  ratio experiendi quotidie in varijs locis datur', Roger Bacon
  (\href{https://doi.org/10.3931/e-rara-16491}{1733}, Pars VI: Scientia
  Experimentalis).}.

\hypertarget{references}{%
\subsection{References}\label{references}}

Ambrosius, A. (12th c.) \emph{Ambrosii Hexaemeron Libri VI}. München:
Bayerische Staatsbibliothek.
\url{https://mdz-nbn-resolving.de/details:bsb00046506}

Augustinus, A. (1961). \emph{Über Den Wortlaut Der Genesis DE GENESI AD
LITTERAM LIBRI DUODECIM: Der Große Genesiskommentar in Zwölf Büchern Zum
Erstenmal in Dt. Sprache von Carl Johann Perl}. Edited by Perl, C. J.
Paderborn: Schöningh.
\url{https://mdz-nbn-resolving.de/details:bsb00046071}

Bacon, R. (1733). \emph{Fratris Rogeri Bacon \ldots{} opus maius
\ldots{} ex MS. codice Dubliniensi}. Edited by Jebb, S. Cum aliis
quibusdam collato, nunc primum edidit. Londini: typis Gulielmi Bowyer.
\url{https://doi.org/10.3931/e-rara-16491}

Jiménez de Cisneros, F. (1517). \emph{Biblia Polyglotta Complutensis}.
Complutum: Arnaldo Guillén de Brocar.
\url{https://doi.org/10.3931/e-rara-46695}

Jung, C. G. (2023). \emph{Collected Works of C. G. Jung: The First
Complete English Edition of the Works of C. G. Jung}. United Kingdom:
Routledge. \url{https://books.google.com/books?id=cX3VEAAAQBAJ}

Mirandola, G., \& Mirandola, G. F. (1601). \emph{Ioannis Pici,
Mirandulae \ldots{} opera quae extant omnia \ldots{}} Basileae: per
Sebastianum Henricpetri. \url{https://doi.org/10.3931/e-rara-61217}

Plato. (1588). \emph{Divini Platonis Opera Omnia}, edited by Ficino, M.,
Conrado, S. Lugduni: apud Joannem Lertout.
\url{https://doi.org/10.3931/e-rara-24349}

Plinius Secundus, G. (1250). \emph{Plinius Secundus Major, Historia
Naturalis. Liber Trigesimus Septimus Manu Recentiore Suppletus Est}.
Latin 6797. Paris: Bibliothèque nationale de France.
\url{https://gallica.bnf.fr/ark:/12148/btv1b550140045/}

Plotinus. (1580). \emph{Plotini Platonicorum facile coryphaei Operum
philosophicorum omnium libri LIV in sex Enneades distributi}. Edited by
Ficinus, M. Basileae: Ad Perneam Lecythum.
\url{https://doi.org/10.3931/e-rara-120889}

\end{document}
